% ============================================================
% RÉSUMÉ
% ============================================================
\newpage
\thispagestyle{plain}

\begin{center}
    {\Large \textbf{Résumé}}\\[1cm]
\end{center}

Ce mémoire détaille la conception architecturale et l'implémentation logicielle d'une plateforme web intelligente, spécifiquement pensée pour la transition numérique de l'entreprise \textit{Assiette Gala Sfaxienne}. Ce système répond à des impératifs opérationnels complexes, visant à centraliser et automatiser les flux métiers : administration du catalogue culinaire, traitement transactionnel des commandes, orchestration des événements, émission automatisée de devis et gouvernance par la donnée.

L'infrastructure applicative retenue s'articule autour d'une architecture orientée microservices. Cette dernière segmente la logique métier en entités autonomes (authentification, gestion du catalogue, commandes, événements et intelligence artificielle), fédérées par une passerelle d'interface (API Gateway). Ce paradigme de conception garantit une haute cohésion, un couplage faible et une scalabilité optimale du dispositif.

Afin de soutenir la prise de décision stratégique, un tableau de bord analytique, enrichi par des composants graphiques dynamiques (Recharts), offre à la direction une visibilité granulaire sur les indicateurs de performance (KPI) et les métriques de vente. En outre, la plateforme intègre des modules métiers avancés, couvrant la gestion logistique des stocks, la relation fournisseurs, la facturation numérisée et l'édition de rapports d'activité.

L'innovation majeure de ce projet réside dans l'intégration d'une couche d'intelligence artificielle contextuelle. Celle-ci se matérialise par un agent conversationnel s'appuyant sur l'approche RAG (Retrieval-Augmented Generation), couplée à une base de données vectorielle FAISS, ainsi que par un moteur de recommandations prédictives. Le cycle de vie du développement a été rigoureusement piloté selon la méthodologie agile Scrum, structurée en cinq itérations (sprints) garantissant une validation fonctionnelle incrémentale.

Les livrables de ce projet prouvent la viabilité technique d'un écosystème numérique souverain, parfaitement aligné avec les contraintes du secteur de la restauration, tout en posant les fondations techniques pour de futures évolutions (déploiement infonuagique, analytique avancée, interface mobile).

\vspace{0.8cm}
\noindent\textbf{Mots-clés :} traiteur, microservices, intelligence artificielle, RAG, FAISS, recommandations, gestion d'événements, tableau de bord, analytics, Scrum.

\clearpage
